\subsection{正算子}

定义 6. --- 设 E 为希尔伯特空间，\(u \in \mathcal{L}(E)\)。我们称 u 是正的并记作 \(u \geqslant 0\)，如果 u 是自伴的且 \(\langle x|u(x) \rangle \geqslant 0\) 对所有 \(x \in E\) 成立。

当 K 等于 C 时，关系式

\[
\langle x | u (x) \rangle \geqslant 0 \quad \text { for   all } \quad x \in E
\]

蕴含 \(u\) 是自伴的（V, p.~2, 注），从而是正的。

用 \(\mathcal{L}_{+}(\mathrm{E})\) 表示 \(\mathcal{L}(\mathrm{E})\) 中所有正元素组成的集；它是实向量空间 \(\mathcal{L}(\mathrm{E})_{[\mathbb{R}]}\)（\(\mathcal{L}(\mathrm{E})\) 的底向量空间）中的一个真尖凸锥。为使 \(u\) 是正的，必要且充分的是：半双线性型 \(\Phi_u\)（在 \(\mathrm{E} \times \mathrm{E}\) 上，与 \(u\) 相伴）是正埃尔米特型。给定 \(u\) 与 \(v\)（都属于 \(\mathcal{L}(\mathrm{E})\)），关系式 \(u - v \geqslant 0\) 也可写作 \(u \geqslant v\) 或 \(v \leqslant u\)；这是 \(\mathcal{L}(\mathrm{E})_{[\mathbb{R}]}\) 上的一个序关系，与它的实向量空间结构相容。

命题 12. --- 设 u 为 \(\mathcal{L}(\mathrm{E})\) 中的自伴（相应地，正的）元素，v 为从 E 到希尔伯特空间 F 中的连续线性映射。则 vuv* 是 \(\mathcal{L}(\mathrm{F})\) 中的自伴（相应地，正的）元素。

因为，我们有 \((vuv^{*})^{*} = v^{**}u^{*}v^{*} = vuv^{*}\)。另一方面，若 \(u \geqslant 0\)，我们有

\[
\langle y | v u v ^ {*} (y) \rangle = \langle v ^ {*} (y) | u (v ^ {*} (y)) \rangle \geqslant 0
\]

对所有 \(y \in \mathbf{F}\) 成立，因而 \(vuv^{*} \geqslant 0\)。

命题 12 特别表明，\(vv^*\) 对所有 \(v \in \mathcal{L}(\mathrm{E};\mathrm{F})\) 都是正的。由于特别地，正交投影算子 \(p\) 满足 \(p = p^2 = pp^*\)，故它是正的。

注. --- 1) 对每个自伴的 \(u\)（属于 \(\mathcal{L}(\mathrm{E})\)），令 \(m(u) = \inf_{||x|| = 1} \langle x|u(x) \rangle\)，\(\mathbf{M}(u) = \sup_{||x|| = 1} \langle x|u(x) \rangle\)。若 E 不恰为 0，则 \(m(u)\) 与 \(\mathbf{M}(u)\) 都是有限的；此外，\(\mathbf{M}(u)\) 是最小的实数 \(\lambda\)，使得 \(u \leqslant \lambda\) . \(1_{\mathrm{E}}\)；而 \(m(u)\) 是最大的实数 \(\mu\)，使得 \(u \geqslant \mu\) . \(1_{\mathrm{E}}\)。显然我们有 \(m(-u) = -\mathbf{M}(u)\) 及 \(\mathbf{M}(-u) = -m(u)\)。显然

\[
\sup \left(| m (u) |, | \mathrm{M} (u) |\right) = \sup _ {\| x \| = 1} \left| \langle x | u (x) \rangle \right|
\]

而命题 9（V, p.~44）蕴含（对 E ≠ \{0\}）

\[
\| u \| = \sup \left(| m (u) |, | \mathbf {M} (u) |\right).\tag{17}
\]

* 当 K 为 C 时，这一公式的另一证明见 TS, I, § 6, No.~8 的命题 14。 * 2) 设 M 与 N 为 E 的两个闭向量子空间，\(p_{\mathbf{M}}\)（相应地 \(p_{\mathbf{N}}\)）为从 E 到 M（相应地 N）上的正交投影算子。则 \(\mathbf{M} \subset \mathbf{N}\) 当且仅当 \(p_{\mathbf{M}} \leqslant p_{\mathbf{N}}\)。因为，我们有 \(p_{\mathbf{M}}^{*}p_{\mathbf{M}} = p_{\mathbf{M}}\)，因而

\[
\left\| p _ {\mathbf {M}} (x) \right\| ^ {2} = \langle p _ {\mathbf {M}} (x) | p _ {\mathbf {M}} (x) \rangle = \langle x | p _ {\mathbf {M}} ^ {*} p _ {\mathbf {M}} (x) \rangle = \langle x | p _ {\mathbf {M}} (x) \rangle
\]

对所有 \(x \in E\) 成立。因此关系式 \(p_{M} \leqslant p_{N}\) 等价于 « \(\|p_{M}(x)\| \leqslant \|p_{N}(x)\|\) 对所有 \(x \in E\) 成立 »。若 \(M \subset N\)，我们有 \(p_{M} = p_{M} p_{N}\)，因而 \(\|p_{M}(x)\| \leqslant \|p_{N}(x)\|\)，因为 \(\|p_{M}\| \leqslant 1\)。反之，若 \(\|p_{M}(x)\| \leqslant \|p_{N}(x)\|\) 对所有 \(x \in E\) 成立，则 \(p_{M}\) 的核包含 \(p_{N}\) 的核，即 \(M^{\circ} \supset N^{\circ}\)，这蕴含 \(M \subset N\)。

命题 13. --- 设 \(\mathcal{H}(\mathrm{E})\) 为希尔伯特空间 E 的所有连续自伴算子组成的集。设 \(\mathcal{F}\) 为 \(\mathcal{H}(\mathrm{E})\) 的一个非空、递增有向且有界的子集。

\begin{enumerate}
\def\labelenumi{(\roman{enumi})}
\tightlist
\item
  集 \(\mathcal{F}\) 有一个上界 \(u_0\)（在 \(\mathcal{H}(\mathrm{E})\) 中）；我们有
\end{enumerate}

\[
\langle x | u _ {0} (x) \rangle = \sup _ {u \in \mathcal {F}} \langle x | u (x) \rangle \quad f o r a l l \quad x \in E.\tag{18}
\]

\begin{enumerate}
\def\labelenumi{(\roman{enumi})}
\setcounter{enumi}{1}
\tightlist
\item
  \(\mathcal{F}\) 的截口滤子收敛于 \(u_0\)（在赋有简单收敛拓扑的空间 \(\mathcal{L}(\mathrm{E})\) 中）。
\end{enumerate}

设 \(\Sigma\) 为 \(\mathcal{F}\) 的截口滤子；对每个 \(u \in \mathcal{H}(\mathrm{E})\)，设 \(\Phi_u\) 为 \(\mathbf{E}\) 上由下式定义的连续埃尔米特型

\[
\Phi_ {u} (x, y) = \langle x | u (y) \rangle .
\]

令

\[
\Psi_ {u} (x) = \Phi_ {u} (x, x)
\]

（其中 \(u \in \mathcal{H}(\mathrm{E})\)，\(x \in \mathrm{E}\)）。由极化公式（V, p.~2），我们有

\[
4 \Phi_ {u} (x, y) = \Psi_ {u} (x + y) - \Psi_ {u} (x - y) \quad \text { if } \quad \mathbf {K} = \mathbf {R}\tag{19}
\]

\[
4 \Phi_ {u} (x, y) = \Psi_ {u} (x + y) - \Psi_ {u} (x - y) - i \Psi_ {u} (x + i y) + i \Psi_ {u} (x - i y) \quad \text { if } \quad \mathbf {K} = \mathbf {C}.\tag{20}
\]

对每个 \(x \in \mathbf{E}\)，映射 \(u \mapsto \Psi_u(x)\)（取值于 \(\mathbf{R}\)）是递增且有界的，因而关于 \(\Sigma\) 有极限。由前面的公式，极限

\[
\lim _ {u, \Sigma} \Phi_ {u} (x, y) = \Phi (x, y)
\]

对 E 的每一对元素 \((x, y)\) 都存在。显然 \(\Phi\) 是 E 上的埃尔米特型。若 \(v_{1} \in F\) 且 \(v_{2}\) 是 F 的一个界，则埃尔米特型 \(f_{1} = \Phi - \Phi_{v_{1}}\) 与 \(f_{2} = \Phi_{v_{2}} - \Phi\) 是正的；存在实数 \(M \geqslant 0\) 使

\[
f _ {1} (x, x) + f _ {2} (x, x) = \Phi_ {v _ {2} - v _ {1}} (x, x) \leqslant \mathbf {M} \| x \| ^ {2},
\]

因而

\[
f _ {1} (x, x) \leqslant \mathbf {M} \| x \| ^ {2}, f _ {2} (x, x) \leqslant \mathbf {M} \| x \| ^ {2} (x \in \mathrm{E});
\]

因而半范数 \(x \mapsto f_i(x, x)^{1/2}\) 在 E 上连续。由于

\[
f _ {2} - f _ {1} = \Phi_ {v _ {2}} + \Phi_ {v _ {1}} - 2 \Phi ,
\]

我们断定 \(x \mapsto \Phi(x, x)\) 是 E 上的连续函数，并由公式 (19) 与 (20)，\(\Phi\) 在 \(E \times E\) 上连续。因此存在（V, p.~16, 推论 2）一个元素 \(u_{0}\)（属于 \(\mathcal{H}(E)\)），使得 \(\Phi = \Phi_{u_{0}}\)。公式 (18) 显然满足，因而 \(u_{0}\) 是 F 在 \(\mathcal{H}(E)\) 中的上界。这就证明了 (i)。

按构造，我们有

\[
\lim _ {u, \Sigma} \langle x | (u _ {0} - u) (x) \rangle = 0 \quad \text { for   all } \quad x \in E.\tag{21}
\]

设 \(v_{1} \in \mathcal{F}\)；给定一个 \(u \in \mathcal{F}\) 满足 \(u \geqslant v_{1}\)，令 \(v = u_{0} - u\)。如果把 Cauchy-Schwarz 不等式应用于 E 上的正埃尔米特型 \(\Phi_v\)，就得到

\[
\begin{array}{r l} \left\| v (x) \right\| ^ {4} & = \left| \Phi_ {v} (v (x), x) \right| ^ {2} \leqslant \Phi_ {v} (v (x), v (x)). \Phi_ {v} (x, x) \\ & = \langle v (x) | v ^ {2} (x) \rangle \langle x | v (x) \rangle \leqslant \| v \| ^ {3} \| x \| ^ {2} \langle x | v (x) \rangle \\ & \leqslant \| u _ {0} - v _ {1} \| ^ {3} \| x \| ^ {2} \langle x | v (x) \rangle , \end{array}
\]

因为由 V, p.~44 的命题 9 有 \(\| v\| \leqslant \| u_0 - v_1\|\)。于是由 (21) 得 \(\lim_{u,\Sigma}\left\| (u_0 - u)(x)\right\| = 0\) 对所有 \(x\in \mathrm{E}\) 成立；这就证明了断言 (ii)。证毕。

特别地，命题 13 可应用于递增且有界的序列 \((u_n)_{n\in \mathbf{N}}\)（元素属于 \(\mathcal{H}(\mathrm{E})\)）的情形。这时存在一个元素 \(v\)（属于 \(\mathcal{H}(\mathrm{E})\)），它由下式刻画

\[
\langle x | v (x) \rangle = \lim _ {n \rightarrow \infty} \langle x | u _ {n} (x) \rangle = \sup _ {n \in \mathbf {N}} \langle x | u _ {n} (x) \rangle \quad (x \in E),
\]

并且有 \(v(x) = \lim_{n\to \infty}u_n(x)\) 对所有 \(x\in E\) 成立。此外，\(v\) 是诸 \(u_{n}\) 组成的集在 \(\mathcal{H}(E)\) 中的上界。

